\section{Comparison with prior multiplication and failure-locus results}
\label{sec:prior-comparison}
The use of secants to account for failure of normal generation is
classical. It is important here to distinguish the fixed complete
series, its varying proper subseries, the Hankel annihilator space,
and the scheme on the parameter space of subseries. The following
comparisons refer to specific results, not to an inference from the
titles of articles.

\begin{center}
\small
\begin{tabular}{@{}p{0.22\textwidth}p{0.34\textwidth}p{0.36\textwidth}@{}}
\toprule
Result compared & Parameter space and conclusion & Use and distinction here\\
\midrule
\cite[Section 1; Theorem 2.1]{CMSV}
& Hankel matrices of indeterminates; determinantal dimensions,
normality and rational singularities.
& These supply the annihilator-rank geometry. They do not identify
the singularities of its varying isotropic-plane image. Our residual
Fitting formulas concern that image.\\[4pt]
\cite{EN}, \cite[Section 2.9]{Kustin}
& A matrix satisfying the expected-grade condition has its
maximal-minor resolution.
& We use this resolution after proving the height. The rank
competition, exceptional incidences and residual normal derivative
are separate arguments.\\[4pt]
\cite[Theorem 0.2; Proposition 2.2]{Ballico96}
& A fixed completely embedded surface; numerical cohomology
conditions force failure on finite linear sections, including a
higher-product version.
& Here the subseries vary, without cross-products. These cited
statements do not give the component or residual Fitting formulas
on $\Gr(c,V_n)^L$.\\[4pt]
\cite[Theorems 1 and 3]{GL86}
& Complete very ample curve series; criteria and secant explanations
for failure of normal generation.
& The ambient complete series on $\Pj^1$ here is normally generated.
The failure comes from varying proper subseries.\\
\bottomrule
\end{tabular}
\end{center}

For the component theorem the additional assertions are the full
rank-incidence optimization, the five exceptional isotropic
incidences, and descent of the signed families. For the global
singularity theorem they are the identification of the polar residual
space and the equality of the restricted cotangent Fitting ideals
with its multiplication minors. The local node theorem further uses
a relative determinant divisor and its normal derivative, not merely
the dimensions of rank strata. The cotangent presentation, Schur
elimination and determinantal resolution themselves are standard
algebraic tools and are not claimed as new constructions.

The closely related paper \cite{Ballico93} requires a distinct
qualification. Its publisher's first page has been inspected: it
introduces higher-order embedding conditions and varying restriction
maps to finite schemes. The remaining pages, containing the theorem
statements needed for a complete comparison, were not available in
the primary text retrieved for this revision. Accordingly no assertion
that its theorems exclude or contain the results here is made. The
comparison with this particular paper remains unverified; the table
above is not a substitute for it. In particular a difference of
parameter-space descriptions on the first page would not by itself
establish a priority claim.

\section{An application to calibrated information recovery}
\label{sec:application-bridge}
Let $J$ be fixed and invertible, let $H(\theta)_{ij}=\theta_{i+j}$,
and set $R_\theta=JH(\theta)J^{\mathsf T}$ in its positive moment
cone. For a prescribed smooth contact, a normal frame $Z$ determines
$U=J^{\mathsf T}\operatorname{ran}Z$. The normal Hessian of the
squared $R_\theta^{-1}$-distance is the inverse of
$Z^{\mathsf T}R_\theta Z$ (Lemma~\ref{lem:normal}). Hence this
contact gives exactly the restriction of the moment functional to
$U^2$. Several independent contacts give the dual of
\eqref{eq:intro-multiplication}. The information-loss locus is
therefore the real positive realization of the failure geometry
studied above, subject to the exact contact hypotheses.

Theorem~\ref{thm:fixed-realization} supplies the positive native
realization and global separation conditions. At a real crossing
with two real branches, Corollary~\ref{cor:node-tail} gives the local
small-singular-value law $\eps^a\log(1/\eps)$. These are applications
of the geometry, not assumptions in its proofs. The appendices give
the complete likelihood, boundary, score-compression, finite-offset
and estimation arguments. In particular they do not identify a
computed feature with an independent noisy distance observation.
